\documentclass[10pt,a4paper]{amsart}
\usepackage[margin=19mm]{geometry}
\usepackage{amsmath,amssymb,mathtools}
\usepackage[scr=rsfs,cal=euler]{mathalfa}
\usepackage{xfrac}
\usepackage[unicode,hidelinks,bookmarks=false]{hyperref}
\newcommand{\ie}{i.e.\ }
\newcommand{\cf}{cf.\ }
\newcommand{\resp}{respectively\ }
\newcommand{\Subsection}{\subsection}
\providecommand{\nobreakdash}{\nobreak\hbox{-}}
\setlength{\emergencystretch}{2em}
\multlinegap0pt
\usepackage{bm}
\usepackage{bbm}
\usepackage{dsfont}
\usepackage{xspace}
\usepackage{cancel}
\usepackage{mathtools}
\usepackage{amsthm}
\makeatletter
\@ifundefined{theorem}{\newtheorem{theorem}{Theorem}}{}
\makeatother
\DeclareMathOperator{\li}{li}
\title[Lines in the prime number graph]{Lines in the prime number graph}
\author{Scott Duke Kominers, Rudi Mrazovi\'c, Carl Pomerance, Patrick Sol\'e}
\date{Source: arXiv:2605.22752v3 (CC BY 4.0)}
\begin{document}
\maketitle

\noindent\emph{Source and licence.} Lines in the prime number graph. Version arXiv:2605.22752v3 (CC BY 4.0), distributed under the Creative Commons Attribution 4.0 International licence, as recorded in the arXiv licence metadata for this exact version and shown on the abstract page https://arxiv.org/abs/2605.22752v3. Full rights notice: licenses/arxiv-2605.22752-CC-BY-4.0.txt.\par\smallskip

\noindent\textit{Editorial scope.} Counted unit: the Theorem whose source label is \texttt{th-L(n)} in the article (source0.tex lines 278--281, source label \texttt{th-L(n)}), delivered together with the article's own complete original proof of it (source0.tex lines 283--355, one \texttt{proof} environment of 73 lines). No mathematics is written, repaired or re-derived: statement and proof are the article's own text line for line. Required context retained uncounted: eq-pnt (lines 257--260). Every definition, display and label of those lines is retained unchanged. Omitted from this package and not redistributed: 1--256, 261--277, 282--282, 356--524. Nothing inside a retained excerpt is omitted or altered: every retained body is delivered exactly as the article's lines are written, so no omission receipt of any kind accompanies this package. Citations: the retained excerpts carry no citation command at all, so no bibliography entry of the article is transcribed and no reference is redistributed. Reference adaptation: none was needed and none was made; every reference inside the delivered unit resolves inside it, so no unresolved reference or citation is produced. Printed slips kept literal: no printed word, symbol, hypothesis or number of the counted statement or of its proof was altered. Layout and class: the article's own class is \texttt{article} with packages including amssymb, latexsym, amsmath, epsfig, amsthm, tikz, color, graphicx, float; the wrapper is the lane's pinned \texttt{amsart} editorial wrapper at 10pt on A4, which restates the theorem-like environments the retained text needs and adds the article's own macro definitions that the retained text needs (\texttt{\textbackslash li}), with extra wrapper packages (bm, bbm, dsfont, xspace, cancel, mathtools). The delivered numbering is the unit's own. Assets: no retained span contains a figure environment, a graphics-inclusion command, a PDF-page inclusion, a table or a TikZ picture; no image file is copied into this package, the package declares no asset dependency and no image file is redistributed with it. Licence: version arXiv:2605.22752v3 of this article is distributed under the Creative Commons Attribution 4.0 International licence, as recorded in the arXiv licence metadata for this exact version and shown on the abstract page https://arxiv.org/abs/2605.22752v3. A full rights notice is delivered at licenses/arxiv-2605.22752-CC-BY-4.0.txt. Characters: every retained excerpt is pure ASCII, so the delivered engine, XeLaTeX with its default Latin Modern text font, reports no missing character.
\begin{equation}
\label{eq-pnt}
|\pi(x)-\li(x)|\le R(x)
\end{equation}

\begin{theorem}
\label{th-L(n)}
We have $L(n)=O(n^{3/4}R(n)^{1/4}/(\log n)^{1/2})$.
\end{theorem}

\begin{proof}
Let $Q$ be a large integer and
consider the Farey sequence of
level $Q$.  Let $a/b, a'/b'$ be two consecutive terms.  Then
\[
a'/b'-a/b = 1/bb',~~b,b'\le Q,~~b+b'>Q,~~\gcd(b,b')=1.
\]
Let $k$ be a large integer, and let
 $u=e^{k+a/b},\,u'=e^{k+a'/b'}$, with $I$ the interval $(u,u']$.  We consider \textbf{inverse prime points} $(p_n,n)$
with $p_n\in I$.   The length $|I|$ of $I$ has
\[
|I|=u'-u=e^{k+a'/b'}-e^{k+a/b}=e^{k+a/b}(e^{1/bb'}-1)=(1+o(1))u/bb', \quad Q\to\infty.
\]
Let $P$ denote the parallelogram bounded by the vertical lines $x=u,\, x=u'$
and the lines with slope $1/\log u=1/(k+a/b)$ through $(u, \li(u)-w)$ and $(u,\li(u)+w)$,
where
\[
w=\frac{|I|^2}{u\log^2u}=(1+o(1))\frac u{(bb')^2\log^2u}.
\]
These lines gain $|I|/\log u$ on the interval $I = (u,u']$.  This is about the same gain as
$\li(x)$ on the interval.  Note that by Taylor's theorem,
\[
\li(u')-\li(u)=\frac{|I|}{\log u}-\left(\frac12+o(1)\right)\frac{|I|^2}{u\log^2 u}=\frac{|I|}{\log u}-\left(\frac12+o(1)\right)w,
\]
and hence the region
\[
    |y-\li(x)|\le w/4,~~x\in I
\]
lies wholly in $P$.
For a given large number $k$ we want to choose $Q$ as large as possible so that the graph of $y=\pi(x)$ for $x\in I$ also lies in $P$.
Since $bb'<Q^2$, we have $w/4 \geq e^k / (8Q^4 k^2)$ for $k$ large, and so by \eqref{eq-pnt} this will be accomplished if we take
\[
Q=\left\lfloor \left(\frac{e^k}{8R(e^{k+1})k^2}\right)^{1/4}\right\rfloor.
\]

We can also do a similar construction by using $u'$ instead of $u$ to determine the slope.
Namely, let $P'$ be the parallelogram bounded by
$x=u,\,x=u'$ and the lines with slope $1/\log u'$ through $(u',\li(u')- w)$ and $(u',\li(u')+ w)$.
With the same choice of $Q$, the prime count $y=\pi(x)$ for $x\in I$ also lies in $P'$ for $k$ large.

Hence all of the inverse prime points $(p_n,n)$ with $p_n\in I$ lie in both $P$ and $P'$.  

If $b<b'$, we
consider those lines with slope $1/\log u=1/(k+a/b)=b/(bk+a)$ that pass through a lattice point in $P$.
Equations of these lines are $(bk+a)y - bx = C$ for different integers $C$,
and so there are at most
\[
(2+o(1))w \cdot (bk+a)\sim 2wb\log u\sim\frac{2u}{bb'^2\log u}\ll\frac{e^k}{bb'^2k}
\]
of them.  If $b>b'$, we take those lines with slope $1/\log u'=b'/(b'k+a')$ which pass through a
lattice point in $P'$; there are $\ll e^k/(b^2b'k)$ of them.
So if we consider all of the lines appearing in this argument for primes in $(e^k,e^{k+1}]$,
the number of them is bounded by a constant times
\begin{equation}
\label{eq-linecount}
\frac{e^k}k\sum\frac{\min\{b,b'\}}{(bb')^2},
\end{equation}
where the sum is over the full Farey dissection of level $Q$.  If $b,b'>Q/2$, then the summand
in \eqref{eq-linecount} is of magnitude $1/Q^3$, and there are fewer than $Q^2$ such pairs,
so the contribution is bounded by $1/Q$.  So, assume $\min\{b,b'\}\le Q/2$.
The contribution to the sum in \eqref{eq-linecount} is at most
\[
2\sum_{b\le Q/2}\frac1b\sum_{Q-b<b'\le Q}\frac1{b'^2}
<2\sum_{b\le Q/2}\frac1b\left(\frac1{Q-b}-\frac1Q\right)
=2\sum_{b\le Q/2}\frac1{(Q-b)Q}\ll\frac1Q.
\]
Thus, the sum in \eqref{eq-linecount} is $O(1/Q)$ and so all of the inverse prime points $(p_n,n)$ with $p_n\in(e^k,e^{k+1}]$ are covered by 
$O(e^k/kQ)$ lines.  With our choice for $Q$
and noting that $R(e^k)\sim R(e^{k+1})$, we have these inverse prime points covered
by $O(e^{3k/4}R(e^k)^{1/4}/k^{1/2})$ lines.  Summing this for $k\le K-1$, we have that the total number of lines that contain
some $(p_n,n)$ for $n\le e^K$ is $O(e^{3K/4}R(e^K)^{1/4}/K^{1/2})$.  Thus, if $n\in(e^{K-1},e^K]$, then
$L(n)=O(n^{3/4}R(n)^{1/4}/(\log n)^{1/2})$.  This completes the proof.
\end{proof}

\end{document}
